
%%%%%%%%%%%%%%%%%%%%%%%%%%%%%%%%%%%%
% Recap/Agenda
%%%%%%%%%%%%%%%%%%%%%%%%%%%%%%%%%%%%
% TODO better formatting
\begin{frame}
    \frametitle{Module 4: Inference}
    \begin{itemize}
        \item \hl{Previously: }Difference of two means (Chapter 7.3)
        \item \hl{This time: }Power calculations for a difference of means (Chapter 7.4)
        \item \hl{Reading: }https://sites.stat.columbia.edu/gelman/research/published/power5r.pdf
        \item \hl{Deadlines/Announcements: }
    \end{itemize}

\end{frame}

%%%%%%%%%%%%%%%%%%%%%%%%%%%%%%%%%%%%
% Learning objectives:
%%%%%%%%%%%%%%%%%%%%%%%%%%%%%%%%%%%%
\begin{frame}
    \frametitle{Learning Objectives}
    \begin{itemize}
        \item \textbf{M4, LO4: Conduct and Interpret t-Tests:} Design, execute, and interpret t-tests for a single population mean, a difference of paired means, and a difference of independent means, calculating the standard error appropriately for each. Describe how to obtain a p-value for a t-test and a critical t-score for a confidence interval.
        \item \textbf{M4, LO6: Calculate Test Power and Evaluate Factors:} Calculate the power of a test for a given effect size and significance level, and explain how the power would change with variations in effect size, sample size, significance level, or standard error.
    \end{itemize}
\end{frame}
